% SC2207 --- Chapter 1: Data Models and ER Design (0->100 ground-up)
\chapter{Data Models, ER Diagrams, and ER-to-Relational Mapping}
\label{ch:models}

\begin{quote}
\emph{``Data knows the truth---if you model it faithfully.''} Before writing a single SQL statement we must decide what entities exist, how they relate, and what constraints govern them. This chapter builds that conceptual foundation from first principles.
\end{quote}

\noindent\fbox{\parbox{\dimexpr\linewidth-2\fboxsep-2\fboxrule}{%
\textbf{From zero: no database assumed.} We start from everyday records (a student list, a library) and build every notion---entity, attribute, relationship, key, abstraction level---from pictures and intuition before any formalism. No prior database course is needed. Later chapters depend only on what is built here.}}
\vspace{0.4em}

\section*{Learning Outcomes}
After this chapter you will be able to:
\begin{enumerate}[label=(L\arabic*)]
    \item distinguish physical, logical, and conceptual abstraction levels;
    \item define and draw Entity--Relationship diagrams with cardinality and participation;
    \item Identify keys, weak entities, and relationship attributes correctly;
    \item map any ER diagram to a relational schema using the standard rules;
    \item critique a design for redundancy and constraint loss.
\end{enumerate}

% ============================================================
\section{Why a Database?}
\label{sec:why-db}

Consider keeping student grades in a spreadsheet shared by email. Soon copies diverge, two tutors overwrite the same cell, and nobody knows which file is authoritative. A file system stores bytes; a \emph{database management system} (DBMS) adds structure, integrity, concurrency, durability, and a declarative query language on top.

\begin{definition}[Database and DBMS]
A \emph{database} is a structured collection of related data modelling some enterprise. A \emph{DBMS} is software that stores, queries, and manages that data while enforcing consistency and supporting concurrent access.
\end{definition}

Think of the DBMS as a disciplined librarian: it catalogues every item, enforces who may borrow what, resolves conflicts when two people request the same book, and guarantees no record is lost even if the lights go out.

Key advantages over flat files: \emph{data independence} (change storage without rewriting apps), \emph{declarative access} (say what you want, not how to fetch it), \emph{transactional guarantees} (Ch.~7), and \emph{centralised control} (security, backup, schema evolution).

\subsection{Levels of Data Abstraction}

\begin{figure}[h]
\centering
\begin{tikzpicture}[scale=0.75, every node/.style={font=\scriptsize},
  box/.style={draw, thick, rounded corners=4pt, align=center, minimum width=3.4cm, minimum height=0.95cm}]
  \node[box, fill=blue!13] (ext) at (0,3.2) {\textbf{External / View level}\\{\tiny user-specific views, e.g.\ student portal}};
  \node[box, fill=green!14] (con) at (0,1.7) {\textbf{Conceptual / Logical level}\\{\tiny ER + relational schema, constraints}};
  \node[box, fill=orange!16] (phy) at (0,0.2) {\textbf{Physical level}\\{\tiny files, indexes, pages on disk (Ch.~5)}};
  \draw[-{Stealth}, thick, blue!60!black] (ext) -- (con) node[midway, right, font=\tiny] {maps};
  \draw[-{Stealth}, thick, blue!60!black] (con) -- (phy) node[midway, right, font=\tiny] {maps};
  \node[align=left, font=\tiny, fill=yellow!10, draw, rounded corners, inner sep=3pt] at (5.2,1.7) {\textbf{Data independence}:\\change physical\\without touching\\conceptual/external};
  \draw[dashed, violet!60!black, thick] (4.0,0.7) -- (4.0,2.7);
\end{tikzpicture}
\caption{Three abstraction levels. Views shield users; conceptual schema is the contract; physical layer optimises storage.}
\label{fig:abstraction}
\end{figure}

Like architectural drawings: the owner's sketch (view), the architect's blueprint (conceptual), and the engineer's construction detail (physical). The blueprint insulates the two.

\subsection{Data Models---A Taxonomy}

A \emph{data model} is the language in which a design is expressed.

\begin{itemize}
    \item \textbf{Hierarchical / Network} (1960s--70s)---trees or graphs of records navigated via pointers (IMS, CODASYL). Fast but rigid; program coupled to traversal path.
    \item \textbf{Relational} (Codd, 1970)---data as tables (relations); queried declaratively by algebra/SQL. Dominant for OLTP.
    \item \textbf{Entity--Relationship}---conceptual picture of entities and relationships; not implemented directly but mapped to relational.
    \item \textbf{Semi-structured / Object models}---JSON, XML, objects (Ch.~8) for flexible or graph-shaped data.
\end{itemize}

This book takes ER $\to$ relational $\to$ SQL $\to$ internals. Start with ER because it thinks in enterprise concepts, not tables.

% ============================================================
\section{The Entity--Relationship Model}
\label{sec:er-model}

An \emph{entity} is a distinguishable object (a particular student). An \emph{entity set} is a collection of similar entities (all students). Each entity has \emph{attributes} (properties like ID, name).

\begin{figure}[h]
\centering
\begin{tikzpicture}[scale=0.76, every node/.style={font=\scriptsize},
  ent/.style={draw, thick, fill=blue!12, minimum width=1.6cm, minimum height=0.75cm, align=center},
  rel/.style={draw, thick, diamond, aspect=1.8, fill=orange!18, inner sep=2pt, align=center},
  attr/.style={draw, thick, ellipse, fill=green!12, inner sep=3pt, align=center}]
  % Left: entity sets
  \node[ent] (stu) at (0,2.2) {Student};
  \node[ent] (crs) at (5.0,2.2) {Course};
  \node[rel] (enr) at (2.5,2.2) {Enrolled};
  \node[attr] (sid) at (-1.3,3.5) {ID};
  \node[attr] (sname) at (0,3.5) {Name};
  \node[attr] (cid) at (5.0,3.5) {Code};
  \node[attr] (ctitle) at (6.5,3.5) {Title};
  \node[attr] (grade) at (2.5,0.85) {Grade};
  \draw[thick] (stu) -- (sid); \draw[thick] (stu) -- (sname);
  \draw[thick] (crs) -- (cid); \draw[thick] (crs) -- (ctitle);
  \draw[thick] (stu) -- (enr) node[midway, above, font=\tiny] {(0..*, 0..*)};
  \draw[thick] (crs) -- (enr);
  \draw[thick] (enr) -- (grade);
  % Underline key
  \draw[thick] (-1.55,3.28) -- (-1.05,3.28);
  \draw[thick] (4.75,3.28) -- (5.25,3.28);
  % Weak entity example on right bottom
  \node[ent, double, fill=violet!10] (sec) at (5.0,0.15) {Section};
  \node[rel, fill=red!12] (has) at (2.5,0.15) {Has};
  \node[ent] (dept) at (0,0.15) {Dept};
  \draw[thick, double] (sec) -- (has); \draw[thick] (dept) -- (has);
  \node[font=\tiny, fill=white, inner sep=1pt] at (2.5,0.55) {identifying};
  \node[font=\tiny, align=center] at (2.5,-0.65) {Double rectangle = weak entity (Section needs Course)};
\end{tikzpicture}
\caption{ER elements: rectangles = entity sets, diamonds = relationships, ovals = attributes, underlining = key; double diamond/rectangle = identifying relationship / weak entity. Grade is a relationship attribute.}
\label{fig:er-elements}
\end{figure}

\subsection{Attributes, Keys, and Domains}

Each attribute has a \emph{domain} (allowed values). Attribute types: \emph{simple} vs.\ \emph{composite} (Address $\to$ Street/City), \emph{single-valued} vs.\ \emph{multi-valued} (double oval), \emph{derived} (dashed oval, e.g.\ age from DOB).

A \emph{superkey} uniquely identifies entities; a \emph{key} is a minimal superkey; a \emph{candidate key} is any minimal key, one chosen as \emph{primary key} (underlined). Weak entities (e.g.\ Section identified only via Course) use a \emph{discriminator} (partial key) plus the owner's key.

\subsection{Relationships, Cardinality, Participation}

A \emph{relationship set} associates entities. \emph{Degree}: binary (most common), ternary, etc. Each entity set's participation has:

\begin{itemize}
    \item \textbf{Cardinality ratio}---1:1, 1:N, N:1, M:N (how many of the other side one entity may relate to). Example: a Course \emph{has} many Sections (1:N).
    \item \textbf{Participation constraint}---\emph{total} (double line: every entity must participate) vs.\ \emph{partial} (single line: optional).
\end{itemize}

An arrow to an entity set means \emph{at most one}; a double line means \emph{at least one}. So a thick arrow-headed double line = exactly one.

\begin{remark}[Attributes on relationships]
Place an attribute on a relationship when it describes the association, not either entity alone (e.g.\ Grade belongs to Enrolled, not to Student or Course). If in doubt, ask: ``Would this value exist without the relationship?''
\end{remark}

% ============================================================
\section{From ER to Relational Schema}
\label{sec:er-to-relational}

Mapping is mechanical once constraints are understood.

\begin{enumerate}
    \item \textbf{Strong entity set} $\to$ table with same attributes; primary key as declared.
    \item \textbf{Weak entity set} $\to$ table with owner's PK + discriminator; PK is composite.
    \item \textbf{Binary 1:N} $\to$ add the ``one'' side's PK as foreign key to the ``many'' side (no separate table). Place relationship attributes on the many side.
    \item \textbf{Binary M:N} $\to$ new join table with PKs of both sides (composite PK); relationship attributes go here.
    \item \textbf{Binary 1:1} $\to$ either side may hold the other's PK (choose the total-participation side to avoid nulls); or a separate table if both partial.
    \item \textbf{Multi-valued attribute} $\to$ separate table (PK of entity + value).
    \item \textbf{N-ary relationship} $\to$ new table with FKs to all participating entity sets.
\end{enumerate}

\begin{figure}[h]
\centering
\begin{tikzpicture}[scale=0.72, every node/.style={font=\tiny}]
  \node[draw, thick, fill=blue!10, minimum width=5.2cm, minimum height=0.85cm, align=left] (s) at (0,2.0) {\textbf{Student}(\underline{ID}, Name)};
  \node[draw, thick, fill=orange!12, minimum width=5.2cm, minimum height=0.85cm, align=left] (c) at (0,0.9) {\textbf{Course}(\underline{Code}, Title)};
  \node[draw, thick, fill=green!12, minimum width=5.8cm, minimum height=0.85cm, align=left] (e) at (0,-0.2) {\textbf{Enrolled}(\underline{StudID, Code}, Grade)~~FKs $\to$ Student, Course};
  \node[draw, thick, fill=violet!10, minimum width=5.8cm, minimum height=0.85cm, align=left] (sec) at (0,-1.3) {\textbf{Section}(\underline{Code, SecNo}, Room)~~FK Code $\to$ Course};
  \draw[-{Stealth}, thick, blue!60!black] (e.north) -- ++(0,0.18) -| (s.south);
  \draw[-{Stealth}, thick, blue!60!black] (e.north) -- ++(0.4,0.18) -| (c.south);
  \node[font=\tiny, fill=yellow!10, draw, rounded corners, inner sep=2pt, align=left] at (5.2,0.35) {Rule of thumb:\\1:N $\to$ FK on N-side\\M:N $\to$ join table\\Weak $\to$ composite PK};
\end{tikzpicture}
\caption{Relational schema resulting from Figure~\ref{fig:er-elements}. Underline = primary key; FKs preserve relationships; null avoidance guides FK placement.}
\label{fig:er-to-rel}
\end{figure}

\begin{example}[Worked mapping]
ER: Student(M)---Enrolled(N)---Course, plus Course 1:N Section (weak). Relational: Student, Course as above; Enrolled join table captures M:N + Grade; Section reuses Course.Code. Cardinalities preserved via keys and foreign-key constraints; total participation becomes \texttt{NOT NULL} on FK.
\end{example}

\begin{remark}[What can be lost?]
ER participation (``every student must enrol in at least one course'') is not enforceable by keys alone---it needs assertions or triggers beyond basic DDL. Flag such constraints when mapping.
\end{remark}

\subsection{Keys in the Relational Target}

After mapping, keys become concrete constraints. A \emph{primary key} (PK) is a minimal set of columns whose values are unique and non-null for every row---exactly the ER key now enforced by the DBMS. A \emph{foreign key} (FK) is a column set whose values must appear as a PK in the referenced table (referential integrity). An FK may be \texttt{NOT NULL} to enforce total participation, or nullable to allow partial participation. This is how the conceptual cardinality survives into tables.

\begin{example}[Keys survive mapping]
In Figure~\ref{fig:er-to-rel}, Enrolled's composite PK $(\mathit{StudID},\mathit{Code})$ prevents a student enrolling twice in the same course. Each column is also an FK, so you cannot enrol a non-existent student. If Course---Section is 1:N total on the Section side, Section.Code is \texttt{NOT NULL} FK.
\end{example}

\subsection{Design Pitfalls to Avoid}

\begin{itemize}
    \item \textbf{Redundancy}---storing Course.Title inside every Enrolled row. Update one copy and miss another: inconsistency.
    \item \textbf{Overloaded keys}---using Name as key when names collide; always prefer stable surrogate or matric number.
    \item \textbf{Misplaced attributes}---putting building capacity on Room when it belongs to Building.
\end{itemize}
A quick sanity check: can you state each attribute as ``X's Y'' with X being the entity/relationship that owns it? If not, ownership is likely wrong.

% ============================================================
\section{Chapter Notes}

The ER model is due to Peter Chen (1976); the abstraction-level view follows ANSI/SPARC (1975). Relational mapping rules are those of Elmasri--Navathe and Silberschatz et al. Modern practice often models first in ER/UML, then generates DDL automatically.

% ============================================================
\section{Exercises}
\label{sec:ch01-ex}
\begin{enumerate}[label=E1.\arabic*, leftmargin=*, itemsep=0.75em]
    \item \textbf{Warm-up: abstraction.} Give a concrete view, conceptual, and physical example for an NTU course-registration system. Why does changing the physical index not require rewriting the student view?
    \item \textbf{ER drawing.} Model NTU's ``Student enrols in Course Offerings taught by Instructor in a Room at a Timeslot.'' Identify entities, keys, the ternary relationship, and where Grade and Capacity belong. Mark cardinalities.
    \item \textbf{Weak entities.} Why is Room not weak but Section is? Design the tables for a weak entity Building $\to$ Room (RoomNo only unique within Building) and state its primary key.
    \item \textbf{Mapping choices.} An ER says Every Course has exactly one Coordinator (Lecturer) but a Lecturer may coordinate zero or one Course. Show two relational mappings (FK on Course vs.\ FK on Lecturer vs.\ separate table) and argue which minimises nulls.
    \item \textbf{Critique.} A novice puts StudentPhone as an attribute of Student (single-valued). Give an instance where this loses information, then redesign using a separate table. Which rule justifies it?
\end{enumerate}
% End of Chapter 1
